\section{Results}
\subsection{CFD Dataset}
\begin{figure}[htpb]
	\centering
	\includegraphics{Graphs/theta_peak_impulse_0.pdf}
	\caption{Contour plot peak impulse distribution overview of PE4 dataset}
	\label{fig:theta_peak_impulse_0}
\end{figure}
The CFD dataset features a cell size of 6.25mm, 100g of PE4 at clear scaled distance of $0.014m/kg^{1/3} \leq z \leq 0.107 m/kg^{1/3}$.
\subsection{Randers-Pehrson Bannister (RPB) Model}
The specific impulse acting on a target is dependent on the angle of incidence of the shock front, theta, with respect to the normal of the target surface. 
From \textcite{randers-pehrsonAirblastLoadingModel1997},
\begin{equation}
	i = i_r cos^2\theta + i_{so}(1 + cos^2\theta - 2cos\theta). 
	\label{eq:RPB}
\end{equation}

Figure \ref{fig:RPB} shows the angle of incidence and slant distance for a point on a target subjected to a spherical air burst, a distance $R$ away from the target centre.
From trigonometry, we have that $\bar{R}^2 = R^2 + x^2 + y^2$, where $x \text{ and } y$ are horizontal and vertical distance from the centre of the target, and $cos \theta = R / \bar{R}$. 
 
\begin{figure}[htpb]
	\centering
	\includegraphics[scale=2, keepaspectratio]{Images/RPB.pdf}
	\caption{Angle of incidence and slant distance for a specific point on a square target}
	\label{fig:RPB}
\end{figure}

The RPB model provides the distribution of impulse as a function of incident angle, theta. 
By inputting into the model the peak reflected specific impulse and the peak reflected specific impulse of the slant distance, $\bar{R}$, we can calculate the specific impulse at any given point on a target.
The specific impulse that is fed into the RPB model is obtained from the MCEER-14-0006 curves (Figure \ref{fig:theta_peak_impulse_MCEER_curves}) by \textcite{shinAirBlastEffectsCivil2014}, which are themselves updated versions of the UFC 3-340-02 blast load curves \parencite{dodUnifiedFacilitiesCriteria2008}. 
\begin{figure}
	\centering
	\includegraphics{Graphs/theta_peak_impulse_MCEER_curves.pdf}
	\caption{MCEER curves for specific impulse}
	\label{fig:theta_peak_impulse_MCEER_curves}
\end{figure}

\begin{figure}[htpb]
	\centering
	\includegraphics{Graphs/theta_peak_impulse_RPB_MCEER_plate.pdf}
	\caption{Contour plot peak impulse distribution of 100g PE4 according to the RPB model with MCEER values. TNT equivalence of PE4 assumed as 1.}
	\label{fig:theta_peak_impulse_RPB_MCEER_plate}
\end{figure}

\subsection{Henrych Model}
From \textcite{henrychDynamicsExplosionIts1979}, we have a theoretical basis to calculate specific impulse distributions. 
The hypothesis of a constant velocity for out flowing particles is introduced:
\begin{equation}
	u = u_x = \text{ constant}.
\end{equation}
For a spherical charge, it is provided that
\begin{equation}
	i = \left(\frac{A_0W}{R^2}\right)cos^2\alpha,
	\label{eq:Henrych_sphere}
\end{equation}
where $R$ is the distance from charge centre to centre of target (shown also in Figure \ref{fig:RPB}) and given
\begin{equation}
	A_0 = \frac{N_{xW} + u_x}{4\pi}, 
\end{equation}
and $N_{xW}$ is the displacement velocity of the outburst surface.

\subsection{Jang and Dharmasena model}
Two alternative modelling procedures have been compared: \cite{jangBlastLoadModel2018} and \cite{dharmasenaMechanicalResponseMetallic2008}.
Both models are derivatives of the RPB model. 
Dharmasena ignored the perpendicular incident impulse and only considered the initial term from Equation \ref{eq:RPB}.
\begin{equation}
i_{cr}(\theta) = i_rexp(-tan^2\theta),
\label{eq:Dharmasena}
\end{equation}
Building from this, \cite{jangBlastLoadModel2018} introduced a weighting function $\beta(z)$ that is a function of scaled distance.
Resulting in the model:
\begin{equation}
i_{cr}(\theta, z) = i_rexp(-\beta(z)tan^2\theta),
\label{eq:Jang}
\end{equation}
where
\begin{equation}
\beta = a(1 + b(z-c))exp(-b(z-c)) + d.
\end{equation}
The coefficients $a$, $b$ and $c$ are determined from finite element simulations of a steel plate blast test whilst $d$ is chosen as 20\% of $\beta$.